%% kapa: exam=1 round=25 year=2014 period=2 form=A subject=appraisal_law subject_ko=부동산관계법규 seq=19 number=19
%% source: Q-Net 감정평가사 기출문제 (cst00309) artlSeq=5101207 file=extracted/1008/5101207/2014년 제25회 감정평가사 1차시험 기출문제/2교시(A형).hwp
%% conversion: kordoc 4.24.51 -> tools/gen1.py (auto); review=proofread-sonnet; flags=none
\begin{question}{19}
부동산 가격공시 및 감정평가에 관한 법령상 표준지 적정가격의 조사ㆍ평가에 관한 설명으로 옳지 않은 것은?
\begin{choices}
  \item 국토교통부장관이 감정평가업자에게 의뢰한 표준지의 적정가격은 감정평가업자가 제출한 조사ㆍ평가액 중 가장 높은 가격을 기준으로 한다.
  \item 국토교통부장관이 표준지의 적정가격을 조사ㆍ평가하고자 할 때에는 둘 이상의 감정평가업자에게 이를 의뢰하여야 한다.
  \item 감정평가업자가 표준지조사평가보고서를 제출하고자 하는 때에는 미리 당해 표준지를 관할하는 시장ㆍ군수 또는 구청장의 의견을 들어야 한다.
  \item 표준지조사평가보고서에는 지역분석조서가 첨부되어야 한다.
  \item 표준지의 적정가격 평가시 표준지에 지상권이 설정되어 있는 때에는 당해 지상권이 존재하지 아니하는 것으로 보고 적정가격을 평가하여야 한다.
\end{choices}
\end{question}
\answer{1}
